\section*{Reading notes}
\subsection*{\texttt{GELFAND-P001475}}
\par\noindent\leavevmode\phantomsection\label{GELFAND-0160-U001}\label{GELFAND-P001475:complex-cancel}\textbf{\texttt{GELFAND-0160-U001}} (complex-cancel).\nopagebreak[4]\par
In parentheses before nabla\_ad squared = 0, a small sign is densely crossed out; a loop and intersecting strokes are visible.\par
Candidates: \textbackslash{}varphi; \textbackslash{}Phi; an illegible crossed-out sign. Unresolved in the diplomatic source.\par
Historical source review: Original pixels x=1995–2110,y=885–985; the original, context at 3× and sign at 10× were viewed. The earlier sign cannot confidently be recovered beneath the cancellations.\par
Current source comparison: The loop remains covered at 5×/10×. Neither candidate field symbol is established.\par\medskip
\par\noindent\leavevmode\phantomsection\label{GELFAND-0160-U002}\label{GELFAND-P001475:morse-function-insert}\textbf{\texttt{GELFAND-0160-U002}} (morse-function-insert).\nopagebreak[4]\par
Above “the fcn. is Morse” are an inserted integral and a small v-shaped sign.\par
Candidates: \textbackslash{}int\textbackslash{}nu; \textbackslash{}int v. Unresolved in the diplomatic source.\par
Historical source review: Original pixels x=580–760,y=1220–1360; viewed at 4× and 10×. The integral is clear; Greek nu and Latin v cannot reliably be distinguished here. A similar sign occurs in (25) on GELFAND-P001477; comparison did not resolve nu/v. The notation was not reconstructed from mathematical expectation.\par
Current source comparison: The integral remains clear, while the following glyph still permits nu or v. No name for the functional is supplied.\par\medskip
\par\noindent\leavevmode\phantomsection\label{GELFAND-0160-U003}\label{GELFAND-P001475:homology-equivalence}\textbf{\texttt{GELFAND-0160-U003}} (homology-equivalence).\nopagebreak[4]\par
After (16), which says any representation of pi\_1, with the addition “or for all local syst. [source abbreviation]”, comes “This is equivalent to”; (17) says M is a homology sphere, i.e. H\_1(X,Z)=0.\par
Candidates: Literally: H\_1(X,\textbackslash{}mathbb Z)=0; Possible notational slip: H\_1(M,\textbackslash{}mathbb Z)=0. Unresolved in the diplomatic source.\par
Historical source review: Region (16)–(17): x=160–1790,y=1480–1995; (16) was viewed at 4× and (17) directly at original resolution. X in (17) was retained. The passage from a condition involving local systems to an ordinary homology sphere requires clarification of hypotheses; this mathematical check was not used to correct the text.\par
Current source comparison: The enlarged views confirm X and integral coefficients. The quantifier over representations is not silently replaced by the adjoint local system associated with one connection; the separate mathematical note states that distinction.\par\medskip
\par\noindent\leavevmode\phantomsection\label{GELFAND-0160-U004}\label{GELFAND-P001475:index-theorem-cancel}\textbf{\texttt{GELFAND-0160-U004}} (index-theorem-cancel).\nopagebreak[4]\par
Between “with” and the abbreviation for “aid”, in the sentence about the usual index theorem, a word is densely crossed out. Only parts of its outer and upper strokes remain visible.\par
Candidates: none recorded. Unresolved in the diplomatic source.\par
Historical source review: Original pixels x=850–1270,y=3040–3210; viewed at 5×. No complete sequence of letters is confidently legible; the cancelled fragment is explicitly marked and the active text preserved.\par
Current source comparison: Two overlapping 5×/10× details still do not yield a complete word.\par\medskip
\subsection*{\texttt{GELFAND-P001476}}
\par\noindent\leavevmode\phantomsection\label{GELFAND-0160-U005}\label{GELFAND-P001476:cotangent-cancel}\textbf{\texttt{GELFAND-0160-U005}} (cotangent-cancel).\nopagebreak[4]\par
Between N= and T*L a sign is crossed out; an oval loop is visible on its left.\par
Candidates: \textbackslash{}Omega; an illegible crossed-out sign. Unresolved in the diplomatic source.\par
Historical source review: Original pixels x=795–940,y=1130–1290; viewed at 3× and 8×. Omega is not to be taken as an established reading.\par
Current source comparison: The oval remnant remains compatible with Omega, but cancellation prevents confirmation.\par\medskip
\par\noindent\leavevmode\phantomsection\label{GELFAND-0160-U006}\label{GELFAND-P001476:below-section-cancel}\textbf{\texttt{GELFAND-0160-U006}} (below-section-cancel).\nopagebreak[4]\par
Below L\_1=L, a short mathematical fragment is crossed out, with a tall initial stroke and a low loop-shaped ending.\par
Candidates: H\_2; L\_2; an illegible crossed-out fragment. Unresolved in the diplomatic source.\par
Historical source review: Original pixels x=1280–1500,y=1295–1470; viewed at 3× and 7×. The earlier letters, possible subscript and cancelling strokes cannot confidently be separated.\par
Current source comparison: The upper and lower strokes remain entangled with cancellation at 10×. Neither the letter nor its subscript is recovered.\par\medskip
\subsection*{\texttt{GELFAND-P001477}}
\par\noindent\leavevmode\phantomsection\label{GELFAND-0160-U201}\label{GELFAND-P001477:one-form-symbol}\textbf{\texttt{GELFAND-0160-U201}} (one-form-symbol).\nopagebreak[4]\par
In (25), on the left inside the angle brackets, and in the following line before “Morse”, is the same v-shaped sign with a tall curved right branch.\par
Candidates: \textbackslash{}nu; v. Unresolved in the diplomatic source.\par
Historical source review: Original coordinates (510,120)–(575,205) and (385,440)–(465,540), both viewed at 10×; overall context (330,25)–(1210,420) at 3×. The body tentatively uses nu; Latin v is not excluded by the shape. Re-examination of GELFAND-P001475 found a similar sign beneath the integral. No name for the functional was added; no dt is visible in (25).\par
Current source comparison: Both instances have the same high curved right branch. They still do not distinguish the alphabets, and no differential is inserted.\par\medskip
\par\noindent\leavevmode\phantomsection\label{GELFAND-0160-U202}\label{GELFAND-P001477:decomposition-letter}\textbf{\texttt{GELFAND-0160-U202}} (decomposition-letter).\nopagebreak[4]\par
After the theorem about M and the phrase about a surface Z comes W=W\_1 union W\_2, boundary W\_i=Z; the first letter in this formula clearly reads W.\par
Candidates: none recorded. Unresolved in the diplomatic source.\par
Historical source review: Original region (330,1115)–(1430,1290) viewed at 3×; theorem context (295,570)–(2150,1130) at 3×. W was retained in the text. It is unclear whether another notation is being introduced or this is a slip; W was not replaced by M.\par
Current source comparison: The first W remains unambiguous at 5×/10×. Its relationship to the preceding M remains unexplained in the literal text.\par\medskip
\par\noindent\leavevmode\phantomsection\label{GELFAND-0160-U203}\label{GELFAND-P001477:cancelled-representation-word}\textbf{\texttt{GELFAND-0160-U203}} (cancelled-representation-word).\nopagebreak[4]\par
Between “Repr. pi\_1(M)” and the dash before “within it” is one word covered by several bold horizontal strokes; a long lower stroke can be distinguished inside the word.\par
Candidates: none recorded. Unresolved in the diplomatic source.\par
Historical source review: Original region (870,1835)–(1090,1970) viewed at 6×; line (300,1800)–(2160,2000) at 3×. Individual letters cannot reliably be separated from cancellation. The fact and position of the cancelled word were retained in the body.\par
Current source comparison: The descender remains visible, but does not establish the cancelled word.\par\medskip
\par\noindent\leavevmode\phantomsection\label{GELFAND-0160-U204}\label{GELFAND-P001477:cylinder-decomposition}\textbf{\texttt{GELFAND-0160-U204}} (cylinder-decomposition).\nopagebreak[4]\par
Drawing (28) shows two convex outer parts joined by two nearly horizontal lines, with arcs at the joins. W\_1, W\_2 and “cylinder” are legible. Short sloping strokes occur at the upper joins.\par
Candidates: none recorded. Unresolved in the diplomatic source.\par
Historical source review: Original region (130,1950)–(1150,2270) viewed at 3×, then compared with the entire 2408×3480 original. The geometric role of the short strokes and faint arcs is not entirely clear, so the drawing was retained through SourceFigure without added directions or connections.\par
Current source comparison: The partial recreation types all three clear labels. It retains gaps at the lower joins and does not interpret the upper diagonal strokes as definite arrowheads.\par\medskip
\par\noindent\leavevmode\phantomsection\label{GELFAND-0160-U205}\label{GELFAND-P001477:cancelled-before-path}\textbf{\texttt{GELFAND-0160-U205}} (cancelled-before-path).\nopagebreak[4]\par
Before the clearly written “Path” is a crossed-out word with a large initial letter resembling Cyrillic Р; its middle and final letters are covered by two horizontal lines.\par
Candidates: none recorded. Unresolved in the diplomatic source.\par
Historical source review: Original region (340,2300)–(560,2420) viewed at 6×; overall lower region (325,2285)–(1980,3100) at 3×. A complete reading of the cancelled word was not established.\par
Current source comparison: The Р-like initial remains only a remnant. The uncancelled “Path” is translated separately.\par\medskip
\par\noindent\leavevmode\phantomsection\label{GELFAND-0160-U206}\label{GELFAND-P001477:four-dimensional-cut}\textbf{\texttt{GELFAND-0160-U206}} (four-dimensional-cut).\nopagebreak[4]\par
The final parenthesis clearly says “M\textasciicircum{}4 is cut into 2 three-dimensional ones”.\par
Candidates: none recorded. Unresolved in the diplomatic source.\par
Historical source review: Original region (325,2285)–(1980,3100) viewed at 3×; the parenthesis itself is approximately (340,2920)–(1700,3070). The dimensional content of the phrase is not clarified in the visible text. All visible words and the exponent 4 were retained, without adding a noun or replacing the dimension.\par
Current source comparison: The words and 4 are clear. A separate dimensional check does not supply the absent noun or repair the phrase.\par\medskip
\subsection*{\texttt{GELFAND-P001478}}
\par\noindent\leavevmode\phantomsection\label{GELFAND-0160-U207}\label{GELFAND-P001478:cancelled-margin-label}\textbf{\texttt{GELFAND-0160-U207}} (cancelled-margin-label).\nopagebreak[4]\par
To the left of “Then there is a determinant” are a smeared item sign and a right parenthesis.\par
Candidates: г) [Cyrillic g]; 1). Unresolved in the diplomatic source.\par
Historical source review: Original region (205,1980)–(325,2120) viewed at 6×; overall context (175,1920)–(2280,2455) at 3×. The lower part of the sign is repeatedly retraced, so the candidate was not accepted as a reading.\par
Current source comparison: The retracing still prevents selecting a letter or numeral.\par\medskip
\par\noindent\leavevmode\phantomsection\label{GELFAND-0160-U208}\label{GELFAND-P001478:cancelled-final-symbol}\textbf{\texttt{GELFAND-0160-U208}} (cancelled-final-symbol).\nopagebreak[4]\par
In the last line before “and so on” is one densely corrected or crossed-out sign with a rounded left part and a lower-right tail.\par
Candidates: a; и [Cyrillic i]. Unresolved in the diplomatic source.\par
Historical source review: Original region (330,3220)–(450,3370) viewed at 6×; overall lower region (290,2420)–(1690,3360) at 3×. The separate sign is accounted for, but no confident reading is available.\par
Current source comparison: The rounded shape remains ambiguous at 10×.\par\medskip
\subsection*{\texttt{GELFAND-P001479}}
\par\noindent\leavevmode\phantomsection\label{GELFAND-0160-U209}\label{GELFAND-P001479:overwritten-section-label}\textbf{\texttt{GELFAND-0160-U209}} (overwritten-section-label).\nopagebreak[4]\par
Before “8 letters are arranged in a cube” is an item sign with two dark rounded parts and a right parenthesis; there are traces of retracing or correction.\par
Candidates: в) [Cyrillic v]; г) [Cyrillic g]. Unresolved in the diplomatic source.\par
Historical source review: Original region (195,320)–(295,430) viewed at 6× and independently re-examined at 10×. Both examinations favor Cyrillic в), but the corrections do not exclude г). The numbering was not adjusted to the expected sequence.\par
Current source comparison: The two rounded parts still favor в), without establishing it. No renumbering is applied.\par\medskip
\par\noindent\leavevmode\phantomsection\label{GELFAND-0160-U210}\label{GELFAND-P001479:cancelled-spin-value}\textbf{\texttt{GELFAND-0160-U210}} (cancelled-spin-value).\nopagebreak[4]\par
After sigma:= and before plus or minus 1 is a short entry densely covered by a series of thick sloping strokes.\par
Candidates: none recorded. Unresolved in the diplomatic source.\par
Historical source review: Original region (2070,1190)–(2225,1320) viewed at 7×; the line and lattice (255,890)–(2375,1590) at 3×. The value of the cancelled entry cannot be recovered; the legible plus or minus 1 was retained separately.\par
Current source comparison: The hatching remains opaque to a reliable value.\par\medskip
\par\noindent\leavevmode\phantomsection\label{GELFAND-0160-U211}\label{GELFAND-P001479:ising-cylinder}\textbf{\texttt{GELFAND-0160-U211}} (ising-cylinder).\nopagebreak[4]\par
The vertical cylinder consists of an upper oval, side lines and a lower arc; several longitudinal lines are inside. At the bottom are a short sloping stroke on the left and a transverse stroke farther right.\par
Candidates: none recorded. Unresolved in the diplomatic source.\par
Historical source review: Original region (270,1880)–(1940,2160) viewed at 3×; the drawing itself is approximately (325,1905)–(535,2150). The connections between the inner lines and the outline, and the role of the short strokes, are unclear. SourceFigure retains the original as the diagram source, without assumptions about orientation or additional edges.\par
Current source comparison: The recreated outline and internal vertical segments preserve visible gaps. The two lower short strokes are traced without supplying arrows or extra edges.\par\medskip
\par\noindent\leavevmode\phantomsection\label{GELFAND-0160-U212}\label{GELFAND-P001479:cancelled-before-eta}\textbf{\texttt{GELFAND-0160-U212}} (cancelled-before-eta).\nopagebreak[4]\par
To the left of eta(tau) is a cancelled entry: the initial angular sign resembles Z, and the following lines are covered by long sloping strokes.\par
Candidates: Z …. Unresolved in the diplomatic source.\par
Historical source review: Original region (355,2230)–(610,2410) viewed at 6×; overall formula block (340,2180)–(2020,2600) at 3×. Possible Z is not accepted as a complete reading of the cancelled entry.\par
Current source comparison: The Z-like remnant still does not establish the cancelled expression.\par\medskip
\par\noindent\leavevmode\phantomsection\label{GELFAND-0160-U213}\label{GELFAND-P001479:eta-q-normalization}\textbf{\texttt{GELFAND-0160-U213}} (eta-q-normalization).\nopagebreak[4]\par
On one line are eta(tau)=q\textasciicircum{}(1/24) product(1−q\textasciicircum{}n) and q=exp(i pi tau); the next has theta\_3(tau)=sum exp(i pi tau n\textasciicircum{}2). No factor 2 is visible in the exponent defining q; the product and sum limits are unwritten.\par
Candidates: none recorded. Unresolved in the diplomatic source.\par
Historical source review: Original formula block (340,2180)–(2020,2600) viewed at 3×; the tau arguments separately in (680,2210)–(895,2530) at 5×. The formulas were left literal. A possible normalization discrepancy was marked without inserting 2, limits, coefficients or other corrections.\par
Current source comparison: The initial exponent details clipped its upper part; corrected 5×/10× views show the complete 1/24. They confirm the literal formula, while the comparison with standard eta and theta conventions is separate.\par\medskip
\par\noindent\leavevmode\phantomsection\label{GELFAND-0160-U214}\label{GELFAND-P001479:last-theta-index}\textbf{\texttt{GELFAND-0160-U214}} (last-theta-index).\nopagebreak[4]\par
In the third modulus in the lower formula, the numerator has theta with a small index most resembling 1; earlier, after “and analog.”, theta\_2 and theta\_4 are clearly written.\par
Candidates: \textbackslash{}theta\_1; \textbackslash{}theta\_4. Unresolved in the diplomatic source.\par
Historical source review: Original region (1280,2870)–(1405,3005) viewed at 10×; an independent re-examination of (1235,2860)–(1415,3010) at 10× confirmed the possibility of 1. The body tentatively retains theta\_1 with an explicit marker. Mathematical expectation of theta\_4 was not used to correct it.\par
Current source comparison: The lower index continues to resemble 1. The competing historical candidate remains recorded rather than replaced from a standard partition function.\par\medskip
\subsection*{\texttt{GELFAND-P001480}}
\par\noindent\leavevmode\phantomsection\label{GELFAND-0160-U401}\label{GELFAND-P001480:top-torus-cycles}\textbf{\texttt{GELFAND-0160-U401}} (top-torus-cycles).\nopagebreak[4]\par
A torus with an inner hole and arcs alpha and beta; the lower arc has a leftward arrowhead, and the right side a downward-rightward arrowhead. Some lines are broken or doubled.\par
Candidates: none recorded. Unresolved in the diplomatic source.\par
Historical source review: Original pixels (250,290)–(560,565). The entire image and a 5× crop were examined. It is uncertain whether the short stroke at the upper edge of the hole is a separate arrow; dotted portions of beta and the outline partly merge. SourceFigure preserves the known labels and directions; the original scheme remains authoritative.\par
Current source comparison: The recreation keeps the two clear arrow directions and labels. The short upper-hole stroke receives no asserted direction; gaps and overlaps are not filled in as hidden connections.\par\medskip
\par\noindent\leavevmode\phantomsection\label{GELFAND-0160-U402}\label{GELFAND-P001480:middle-cut-torus}\textbf{\texttt{GELFAND-0160-U402}} (middle-cut-torus).\nopagebreak[4]\par
A cut torus with two oval ends, beta above and theta by the cut; a dense angular stroke lies at the left inner edge.\par
Candidates: none recorded. Unresolved in the diplomatic source.\par
Historical source review: Original pixels (275,1365)–(610,1605); overall region (220,1210)–(2060,1760), enlarged 3×. The cut outlines are visible, but the inner stroke cannot confidently be separated from an arrowhead. No gluing direction was added.\par
Current source comparison: The two oval ends and clear beta and theta are traced. The left angular remnant is retained as a stroke rather than given an arrow direction.\par\medskip
\par\noindent\leavevmode\phantomsection\label{GELFAND-0160-U405}\label{GELFAND-P001480:bottom-point-pair}\textbf{\texttt{GELFAND-0160-U405}} (bottom-point-pair).\nopagebreak[4]\par
The first index in the pair below the last formula is crowded beneath z; j, a horizontal stroke and 1 are visible, with a possible 2 before j. The second index reads 2j.\par
Candidates: (z\_\{2j-1\},z\_\{2j\}); (z\_\{j-1\},z\_\{2j\}); (z\_\{2j+1\},z\_\{2j\}). Unresolved in the diplomatic source.\par
Historical source review: Original pixels (260,2970)–(610,3170). The entire lower block at 3× and a separate 5× crop were examined. The index was not reconstructed from expected paired numbering of the handles.\par
Current source comparison: The additional enlargements do not reliably distinguish all strokes in the first index. The second 2j remains clear.\par\medskip
\par\noindent\leavevmode\phantomsection\label{GELFAND-0160-U406}\label{GELFAND-P001480:bottom-cancelled-word}\textbf{\texttt{GELFAND-0160-U406}} (bottom-cancelled-word).\nopagebreak[4]\par
After “to which” a word is densely crossed out; next comes “attach the (thin) handles”.\par
Candidates: none recorded. Unresolved in the diplomatic source.\par
Historical source review: Original pixels (1070,3000)–(1340,3170), context (1070,3000)–(1660,3170). Viewed at 3× and 5×. Loops and a long final stroke can be distinguished, but the earlier word beneath cancellation was not recovered.\par
Current source comparison: The surviving loops still do not form a secure complete word.\par\medskip
\subsection*{\texttt{GELFAND-P001481}}
\par\noindent\leavevmode\phantomsection\label{GELFAND-0160-U407}\label{GELFAND-P001481:top-cft-sets}\textbf{\texttt{GELFAND-0160-U407}} (top-cft-sets).\nopagebreak[4]\par
An oval CFT’s, an inner oval RCFT’s, and a line from ISING to a small cross or arrowhead inside the inner oval.\par
Candidates: a small cross marking ISING; the head of a pointer arrow from ISING. Unresolved in the diplomatic source.\par
Historical source review: Original pixels (330,70)–(840,505). The whole frame and crop (290,60)–(1350,640) were examined at 3×. The ovals and words are legible; the pointer ending around (570,310) consists of several overlaid strokes.\par
Current source comparison: All three words are typed in the recreation. The dark endpoint remains a short crossing mark without a claimed arrowhead or separately defined node.\par\medskip
\par\noindent\leavevmode\phantomsection\label{GELFAND-0160-U408}\label{GELFAND-P001481:rational-cancelled-prefix}\textbf{\texttt{GELFAND-0160-U408}} (rational-cancelled-prefix).\nopagebreak[4]\par
Before RATIONAL, I and several densely cancelled letters following it are visible on the left.\par
Candidates: I [illegible crossed-out letters]. Unresolved in the diplomatic source.\par
Historical source review: Original pixels (300,500)–(490,610). Viewed at 3× and 7×. The initial I was retained separately; the inner letters of the cancelled fragment were neither read nor supplied from the topic.\par
Current source comparison: The initial I remains distinct from the unreadable letters. The cancelled prefix is not expanded to IRRATIONAL.\par\medskip
\par\noindent\leavevmode\phantomsection\label{GELFAND-0160-U410}\label{GELFAND-P001481:torus-period-convention}\textbf{\texttt{GELFAND-0160-U410}} (torus-period-convention).\nopagebreak[4]\par
Both handwritten lines about F\_j have tau+2 pi i; above is theta→theta+2 pi, and at the bottom of the page T:tau→tau+1.\par
Candidates: the period formulas with tau+2 pi i as written; T:tau→tau+1 as written below. Unresolved in the diplomatic source.\par
Historical source review: Original pixels (270,770)–(1570,1320) and (950,2950)–(1440,3200). Viewed at 3×. A possible change of convention or source error in the period of tau was noted. All three visible entries were retained without correction.\par
Current source comparison: An initial first-period detail was blank and the initial T detail actually showed the S row. Corrected focused 5×/10× views confirm both tau+2 pi i and T:tau→tau+1. Their compatibility is discussed only under an explicitly stated rescaling.\par\medskip
\par\noindent\leavevmode\phantomsection\label{GELFAND-0160-U411}\label{GELFAND-P001481:monodromy-cancelled-word}\textbf{\texttt{GELFAND-0160-U411}} (monodromy-cancelled-word).\nopagebreak[4]\par
After “M = monodromy matrix.” is a small densely crossed-out word.\par
Candidates: none recorded. Unresolved in the diplomatic source.\par
Historical source review: Original pixels (1310,1270)–(1520,1350). Context at 3× and a crop at 8× were examined. Initial loops could belong to several letters; no reliable complete candidate is available.\par
Current source comparison: The loops remain interrupted by cancellation.\par\medskip
\par\noindent\leavevmode\phantomsection\label{GELFAND-0160-U412}\label{GELFAND-P001481:ising-product-conventions}\textbf{\texttt{GELFAND-0160-U412}} (ising-product-conventions).\nopagebreak[4]\par
The source writes sqrt(theta\_2/eta)=q\textasciicircum{}(1/8) product(1+q\textasciicircum{}n), sqrt(theta\_3/eta)=product(1−q\textasciicircum{}(n−1/2)), sqrt(theta\_4/eta)=product(1+q\textasciicircum{}(n−1/2)), q=exp(2 pi i tau); there are no lower product bounds or further factors.\par
Candidates: none recorded. Unresolved in the diplomatic source.\par
Historical source review: Original pixels (310,1660)–(1600,2200); exponent separately (665,1670)–(805,1790). Viewed at 3× and 8×. The exponent 1/8, signs and absence of further normalization factors were reproduced literally. Possible errors or a different convention for theta and eta were not corrected; expected formulas were not substituted.\par
Current source comparison: The 5×/10× details confirm 1/8, minus in the theta\_3 line and plus in the theta\_4 line. Standard product identities are derived in a separate mathematical note, with the nome and root convention stated.\par\medskip
\subsection*{\texttt{GELFAND-P001482}}
\par\noindent\leavevmode\phantomsection\label{GELFAND-0160-U414}\label{GELFAND-P001482:top-genus-two-twists}\textbf{\texttt{GELFAND-0160-U414}} (top-genus-two-twists).\nopagebreak[4]\par
A two-handled surface labeled alpha\_1,beta\_1,alpha\_2,beta\_2,alpha\_3; loops and segments cross or overlap the hole outlines.\par
Candidates: none recorded. Unresolved in the diplomatic source.\par
Historical source review: Original pixels (330,190)–(900,375), context (140,30)–(1050,535), enlarged 3×. The labels were preserved, but all hidden and visible portions of the curves and their joins cannot confidently be recovered. No further crossings or arrows were introduced.\par
Current source comparison: The partial recreation keeps the five labels and visible curve fragments. It leaves gaps where overlaps obscure continuity; homology calculations do not determine the missing drawing topology.\par\medskip
\par\noindent\leavevmode\phantomsection\label{GELFAND-0160-U415}\label{GELFAND-P001482:ln-integration-symbol}\textbf{\texttt{GELFAND-0160-U415}} (ln-integration-symbol).\nopagebreak[4]\par
In two lines with L\_n the same variable resembles zeta or xi. A small c/C is below the first contour integral. No differential is visible.\par
Candidates: variable zeta, contour c; variable xi, contour c; variable zeta, contour C; variable xi, contour C. Unresolved in the diplomatic source.\par
Historical source review: Original pixels (410,2330)–(1850,2680); the first integral separately (695,2390)–(900,2540). Viewed at 3× and 6×. The main reading types zeta and c with an explicit alternatives marker; the differential absent from the source was not added.\par
Current source comparison: The contour label still permits c/C and the variable zeta/xi. No differential or residue normalization is inserted.\par\medskip
\par\noindent\leavevmode\phantomsection\label{GELFAND-0160-U416}\label{GELFAND-P001482:rho-trailing-fragment}\textbf{\texttt{GELFAND-0160-U416}} (rho-trailing-fragment).\nopagebreak[4]\par
After rho\_i:A→End H\_i and a comma, another rho\_i and a rounded sign resembling C or O are written. No continuation is visible to the right.\par
Candidates: rho\_i C; rho\_i O. Unresolved in the diplomatic source.\par
Historical source review: Original pixels (1100,2740)–(1390,2855), context (360,2690)–(1420,2950). Viewed at 3× and 6×. This is a short visible source fragment, rather than an omitted continuation; its meaning was not supplied. A separating row of small strokes begins beneath it.\par
Current source comparison: The short fragment and the separator remain distinct. Its unfinished meaning is retained.\par\medskip
\subsection*{\texttt{GELFAND-P001483}}
\par\noindent\leavevmode\phantomsection\label{GELFAND-0160-U601}\label{GELFAND-P001483:planck-unit}\textbf{\texttt{GELFAND-0160-U601}} (planck-unit).\nopagebreak[4]\par
After “the Planck length is”, the entry is approximately 10\textasciicircum{}(-33) m; no separate Cyrillic s is visible before m.\par
Candidates: approximately 10\textasciicircum{}(-33) m. Unresolved in the diplomatic source.\par
Historical source review: Original 2408×3480 pixels, region (1640,205,1875,285), enlarged 8×; also checked in the whole frame and text strip at 3×. The unit m is confidently legible and was retained. Possible source error: the exponent corresponds to the common length notation in centimeters, but the manuscript unit looks like m; no correction to cm was made. What remains uncertain is the author’s intention, not the presence of a separately visible Cyrillic s.\par
Current source comparison: The visible unit is m. The separate comparison with the SI Planck length does not change the diplomatic exponent or unit.\par\medskip
\par\noindent\leavevmode\phantomsection\label{GELFAND-0160-U602}\label{GELFAND-P001483:sk-annotation}\textbf{\texttt{GELFAND-0160-U602}} (sk-annotation).\nopagebreak[4]\par
A short line from S\_K leads to a three-line note: the first line is “compact”, the second resembles “конд. на”, and the third ends “пр-ва”. The beginning of the third line is joined.\par
Candidates: конд. на утолщ. пр-ва [source abbreviations; not expanded]; конд. на уплотн. пр-ва [source abbreviations; not expanded]; конд. на уполщ. пр-ва [source abbreviations; not expanded]; конд. на уплощ. пр-ва [source abbreviations; not expanded]. Unresolved in the diplomatic source.\par
Historical source review: Original coordinates (1880,955,2280,1230), enlarged 5×, then the entire frame again. An independent second examination confirmed “compact” and favored “конд. на” over “конф. на”, but did not resolve the third line. The abbreviations were not expanded.\par
Current source comparison: The initial third-line detail was mostly blank. The corrected focused 5×/10× details cover the joined beginning, which remains unreadable enough to distinguish the four historical candidates. No compactification phrase is supplied.\par\medskip
\par\noindent\leavevmode\phantomsection\label{GELFAND-0160-U603}\label{GELFAND-P001483:cft-cancellation}\textbf{\texttt{GELFAND-0160-U603}} (cft-cancellation).\nopagebreak[4]\par
Between “Therefore” and “this is interesting”, one word is crossed out with several horizontal strokes; its visible remnants permit “имеет” or “имеют”.\par
Candidates: имеет [has; singular]; имеют [have; plural]. Unresolved in the diplomatic source.\par
Historical source review: Original coordinates (510,1940,740,2040), enlarged 7×; the complete text was checked at 3×. Covered letter strokes cannot be separated from the cancelling line; cancellation was retained in editorial square brackets.\par
Current source comparison: The obscured ending remains unresolved. Neither grammatical number is silently supplied.\par\medskip
\par\noindent\leavevmode\phantomsection\label{GELFAND-0160-U604}\label{GELFAND-P001483:field-overwrite}\textbf{\texttt{GELFAND-0160-U604}} (field-overwrite).\nopagebreak[4]\par
In the line “x\textasciicircum{}beta, … (x) is a field”, the field sign before the parenthesis is densely overwritten; a vertical and overlapping loops are visible.\par
Candidates: \textbackslash{}varphi(x); \textbackslash{}Phi(x). Unresolved in the diplomatic source.\par
Historical source review: Original coordinates (1000,2570,1190,2680), enlarged 8×, additionally the entire lower part at 3×. The body gives a tentative varphi reading with a linked Unclear marker; the reading is not asserted. The superscript beta on x was separately checked at 8× in (820,2560,940,2670) and confidently read.\par
Current source comparison: The overwritten field letter still permits both candidates. Beta is clear independently of that overwrite.\par\medskip
\par\noindent\leavevmode\phantomsection\label{GELFAND-0160-U605}\label{GELFAND-P001483:contour-ending}\textbf{\texttt{GELFAND-0160-U605}} (contour-ending).\nopagebreak[4]\par
After “Phi\{C\} is a fcn. of”, the stem “контур” is visible with a densely overwritten ending.\par
Candidates: контура [of a contour]; контуров [of contours]. Unresolved in the diplomatic source.\par
Historical source review: Original coordinates (770,2820,1040,2950), enlarged 6×, overall lower fragment at 3×. Number and case of the ending were not reconstructed from mathematical context; the visible ending is accounted for by the marker.\par
Current source comparison: The ending remains ambiguous. The English body explicitly records the singular/plural alternative.\par\medskip
\par\noindent\leavevmode\phantomsection\label{GELFAND-0160-U606}\label{GELFAND-P001483:eq1-first-symbol}\textbf{\texttt{GELFAND-0160-U606}} (eq1-first-symbol).\nopagebreak[4]\par
The first field sign inside the angle brackets of (1) is overwritten; the other comparable field signs read varphi.\par
Candidates: \textbackslash{}varphi(x\_1); \textbackslash{}Phi(x\_1). Unresolved in the diplomatic source.\par
Historical source review: Original coordinates (170,2955,740,3105), enlarged 5×, also the whole frame. The first index x\_1 was separately checked by comparison with the right-hand part (1190,2965,1660,3105), at 5×. The positive exponent exp(Sigma(phi)) and absence of a normalization factor were retained without additions.\par
Current source comparison: The complete details show x\_1 on the left and x\_n at the right. An initial third detail duplicated a right-hand factor and is not evidence for x\_1; the first two correctly located details cover the required indices. The first field glyph remains uncertain.\par\medskip
\par\noindent\leavevmode\phantomsection\label{GELFAND-0160-U607}\label{GELFAND-P001483:particle-diagram}\textbf{\texttt{GELFAND-0160-U607}} (particle-diagram).\nopagebreak[4]\par
At the lower left are three lines labeled x\_1,x\_2,x\_3: the upper one is wavy, and the other two rise toward it to the right. Two places where lines approach have dark strokes or points; the rightmost stroke also resembles an arrowhead.\par
Candidates: two joining points of the three incoming lines; one joining point and a rightward arrowhead. Unresolved in the diplomatic source.\par
Historical source review: Original coordinates (180,3105,660,3390), examined within (100,2515,2230,3410) at 3× and again in the whole frame. Line continuity and the character of the rightmost stroke were not established. The diagram is fully accounted for through SourceFigure, labels and neighboring text are preserved, and topology was not completed by invention.\par
Current source comparison: The 10× details show both upper approaches. An initial lower detail was mostly blank; corrected lower-label details confirm x\_3. The recreation types all labels and leaves small gaps at the uncertain upper contacts, with no asserted node or arrowhead.\par\medskip
\subsection*{\texttt{GELFAND-P001484}}
\par\noindent\leavevmode\phantomsection\label{GELFAND-0160-U608}\label{GELFAND-P001484:eq2-relation}\textbf{\texttt{GELFAND-0160-U608}} (eq2-relation).\nopagebreak[4]\par
Between the closing angle bracket and integral in (2) is a densely overwritten sign; a clear pointed end faces right, while the central horizontal strokes are fused.\par
Candidates: \textbackslash{}longrightarrow; \textbackslash{}Rightarrow; =. Unresolved in the diplomatic source.\par
Historical source review: Original coordinates (827,180,965,325), enlarged 8×, also the upper text strip at 3× and the whole frame. A rightward arrow over an earlier sign is more likely, but the arrow type and earlier sign are not established; equality was not substituted by analogy with (1).\par
Current source comparison: The pointed end remains visible, but the horizontal layers do not settle the relation.\par\medskip
\par\noindent\leavevmode\phantomsection\label{GELFAND-0160-U609}\label{GELFAND-P001484:eq5-exponent-overwrite}\textbf{\texttt{GELFAND-0160-U609}} (eq5-exponent-overwrite).\nopagebreak[4]\par
In the exponential exponent in (5), a minus, an overwritten letter resembling m, and L are visible; further strokes lie above and through m.\par
Candidates: -mL. Unresolved in the diplomatic source.\par
Historical source review: Original coordinates (1080,1260,1215,1340), enlarged 10×; middle part at 3×. The final shape resembles m, but the earlier letter or nature of the correction is illegible. The body marks the overwrite; no coefficients or factors were reconstructed.\par
Current source comparison: The visible candidate remains −mL. The earlier content of the overwritten letter is not recovered.\par\medskip
\par\noindent\leavevmode\phantomsection\label{GELFAND-0160-U610}\label{GELFAND-P001484:eq6-measure}\textbf{\texttt{GELFAND-0160-U610}} (eq6-measure).\nopagebreak[4]\par
After (partial phi)\textasciicircum{}2, the first measure in (6) has a sign with a closed lower loop and an arc above: partial\textasciicircum{}(26)x. In the second integral, d\textasciicircum{}(26)x has a clearly different shape.\par
Candidates: \textbackslash{}partial\textasciicircum{}\{26\}x. Unresolved in the diplomatic source.\par
Historical source review: Original coordinates (997,1945,1183,2063), enlarged 8×, overall fragment at 3×. Partial is confidently read and was retained without replacement by the expected differential d. A possible source notation error is recorded here; the author’s meaning was not corrected.\par
Current source comparison: The two different measure glyphs remain clear. Conditional variational calculations explicitly assume an ordinary volume form instead of redefining the first handwritten sign.\par\medskip
\par\noindent\leavevmode\phantomsection\label{GELFAND-0160-U611}\label{GELFAND-P001484:eq6-cubic-free}\textbf{\texttt{GELFAND-0160-U611}} (eq6-cubic-free).\nopagebreak[4]\par
In (6), the second summand is unambiguously plus integral phi cubed d\textasciicircum{}(26)x; before the formula, the action is said to “describe a free particle” [source abbreviations].\par
Candidates: +\textbackslash{}int\textbackslash{}varphi\textasciicircum{}3d\textasciicircum{}\{26\}x. Unresolved in the diplomatic source.\par
Historical source review: Original coordinates of the small symbol (1302,1950,1406,2077), enlarged 8×; the formula and introductory line were checked in (100,900,1710,2110) at 3×. The cubic term disagrees with the usual meaning of a free action; it is not crossed out in the source and was neither removed nor corrected. A possible mathematical error by the author remains unresolved.\par
Current source comparison: The cubed exponent and uncancelled plus sign are clear. The “free particle” wording was also checked in the whole original; an isolated formula detail is not treated as evidence for that introductory wording. The separate third-variation check does not alter either.\par\medskip
\par\noindent\leavevmode\phantomsection\label{GELFAND-0160-U612}\label{GELFAND-P001484:vertex-index}\textbf{\texttt{GELFAND-0160-U612}} (vertex-index).\nopagebreak[4]\par
In “V\_N… are vertex operators”, a small loop-shaped sign with a tall right vertical sits beside N below the line; a separate dot cannot be distinguished.\par
Candidates: V\_\{N\_j\}; V\_\{N\_d\}. Unresolved in the diplomatic source.\par
Historical source review: Original coordinates (400,2570,584,2710), enlarged 7×, lower part at 3×. Context permits j, but the handwriting resembles d; the tentative j in the body is linked to an Unclear marker.\par
Current source comparison: The fresh 5×/10× views show a tiny upper dash or dot-like stroke. This qualifies the historical observation about an indistinguishable dot, without overwriting it; the small mark still does not decide j versus d.\par\medskip
\par\noindent\leavevmode\phantomsection\label{GELFAND-0160-U613}\label{GELFAND-P001484:branching-diagram}\textbf{\texttt{GELFAND-0160-U613}} (branching-diagram).\nopagebreak[4]\par
Drawing (8): a nearly vertical left edge and curved upper and lower boundaries lead to an inner nearly vertical line. Four curves continue rightward from there, with the two inner ones crossing. The word “branching” and drawing are connected by a large arrow with a leftward head.\par
Candidates: a crossing of inner curves without an additional node; a connection of the inner curves at the crossing. Unresolved in the diplomatic source.\par
Historical source review: Original coordinates of drawing and arrow (600,2715,1480,2980), examined within (95,2710,2190,3420) at 3× and in the entire frame. Stroke geometry was described, but their connections at the inner vertical and the meaning of the crossing were not determined. SourceFigure preserves the diagram as a source; no new vertices or links were added.\par
Current source comparison: Two initial lower details mainly showed the pointer or blank area. Corrected 5×/10× views cover the crossing and bottom joins: no explicit node is established, and a lower gap remains near the inner vertical. The recreation keeps those gaps and the clear large leftward pointer.\par\medskip
\par\noindent\leavevmode\phantomsection\label{GELFAND-0160-U614}\label{GELFAND-P001484:branching-cancellation}\textbf{\texttt{GELFAND-0160-U614}} (branching-cancellation).\nopagebreak[4]\par
After “String theory:” and before “forbidden”, a short fragment is crossed out with several strokes; a lower loop of its first letter and a tall final stroke remain.\par
Candidates: зап [cancelled word beginning]; запр [cancelled word beginning]. Unresolved in the diplomatic source.\par
Historical source review: Original coordinates (670,2930,815,3040), enlarged 8×, overall lower strip at 3×. It resembles an abandoned beginning of the following word, but the precise cancelled letters are unreadable. Cancellation is represented in square brackets, without restoring the full word “forbidden”.\par
Current source comparison: The visible remnants still support only uncertain short beginnings, not a complete cancelled word.\par\medskip
